\documentclass[UTF8,fontset=none,zihao=5,twoside,openany]{ctexbook}
\usepackage{bookstyle}
\usepackage{float}
\let\oldfigure\figure \let\endoldfigure\endfigure
\renewenvironment{figure}[1][]{\oldfigure[H]}{\endoldfigure}
\renewcommand{\cite}[1]{}
\setcounter{chapter}{4}
\pagestyle{plain}
% 变体 2：只换字体——图中文字改用正文字体（中文宋体、西文 Times），字号两级 9 / 7.5 pt
\newcommand{\fontbook}{%
  \renewcommand{\figfont}{\rmfamily}%
  \renewcommand{\figsizes}{\def\baselinestretch{1}%
    \renewcommand{\normalsize}{\fontsize{9}{12}\selectfont}%
    \renewcommand{\small}{\fontsize{9}{12}\selectfont}%
    \renewcommand{\footnotesize}{\fontsize{7.5}{10}\selectfont}%
    \renewcommand{\scriptsize}{\fontsize{7.5}{10}\selectfont}%
    \renewcommand{\tiny}{\fontsize{7.5}{10}\selectfont}}}
% 变体 3：再去掉颜色与圆角——只用黑、灰两种墨色
\newcommand{\monobook}{%
  \colorlet{accent}{black}\colorlet{accentpale}{black!12}\colorlet{warn}{black}%
  \colorlet{figfill}{black!6}\colorlet{figline}{black}\colorlet{rulegray}{black!55}\colorlet{ink}{black}}
\begin{document}

\section*{样张三：“AI 味”从哪里来——三个版本对比}
同一张图做三个版本，每次只改一处，便于判断是哪一个因素在起作用：
\begin{enumerate}
  \item \textbf{现行版}：微软雅黑、青绿强调色、圆角框（已修正间距）；
  \item \textbf{只换字体}：图中文字改为与正文相同的宋体 + Times，其余不变；
  \item \textbf{字体 + 配色 + 直角}：在 2 的基础上去掉彩色，只用黑与灰，框改为直角。
\end{enumerate}

\Needspace{70mm}\setcounter{figure}{2}\noindent\textbf{图 4-3\quad 版本 1：现行版}
\begin{figure}[htbp]
\centering
\begin{tikzpicture}[fig, x=1mm, y=1mm]
  % CPU：少量复杂核心 + 大缓存
  \begin{scope}[shift={(-48,0)}]
    \draw[figline, fill=white] (-17,-14) rectangle (17,14);
    \foreach \x/\y in {-8/7,8/7,-8/-1,8/-1}
      \draw[fill=figfill] (\x-6,\y-3.5) rectangle (\x+6,\y+3.5);
    \foreach \x/\y in {-8/7,8/7,-8/-1,8/-1} \node[font=\figfont\scriptsize] at (\x,\y) {核心};
    \draw[fill=white, draw=rulegray] (-14,-12) rectangle (14,-6);
    \node[font=\figfont\scriptsize] at (0,-9) {大容量缓存};
    \node[anchor=north] at (0,-16) {CPU};
    \node[note, anchor=north] at (0,-21) {少量复杂核心};
  \end{scope}
  % GPU：大量简单核心
  \begin{scope}[shift={(0,0)}]
    \draw[figline, fill=white] (-17,-14) rectangle (17,14);
    \foreach \i in {0,...,7} \foreach \j in {0,...,5}
      \draw[fill=figfill, line width=0.3pt] (-14.8+\i*3.8,-11+\j*3.8) rectangle ++(3,3);
    \node[anchor=north] at (0,-16) {GPU};
    \node[note, anchor=north] at (0,-21) {大量简单并行单元};
  \end{scope}
  % NPU：矩阵运算阵列
  \begin{scope}[shift={(48,0)}]
    \draw[figline, fill=white] (-17,-14) rectangle (17,14);
    \foreach \i in {0,...,5} \foreach \j in {0,...,5}
      \draw[fill=accentpale, draw=accent, line width=0.3pt] (-12+\i*4.2,-12+\j*4.2) rectangle ++(3.4,3.4);
    \foreach \j in {0,...,5} \draw[->, accent, line width=0.4pt] (-16,-10.3+\j*4.2) -- (-12.4,-10.3+\j*4.2);
    \node[anchor=north] at (0,-16) {NPU / ASIC};
    \node[note, anchor=north] at (0,-21) {专用矩阵运算阵列};
  \end{scope}
  % 光谱箭头
  \draw[<->, line width=0.7pt] (-65,-30) -- (65,-30);
  \node[note, anchor=north west] at (-65,-31) {通用性、灵活性更强};
  \node[note, anchor=north east] at (65,-31) {专用计算效率更高};
\end{tikzpicture}
\caption{三类计算芯片的结构取舍（示意）。从左到右，芯片的通用性逐渐降低，矩阵运算的吞吐与能效逐渐提高}
\label{fig:chips}
\end{figure}

\Needspace{70mm}\setcounter{figure}{2}\noindent\textbf{图 4-3\quad 版本 2：只换字体}
{\fontbook\begin{figure}[htbp]
\centering
\begin{tikzpicture}[fig, x=1mm, y=1mm]
  % CPU：少量复杂核心 + 大缓存
  \begin{scope}[shift={(-48,0)}]
    \draw[figline, fill=white] (-17,-14) rectangle (17,14);
    \foreach \x/\y in {-8/7,8/7,-8/-1,8/-1}
      \draw[fill=figfill] (\x-6,\y-3.5) rectangle (\x+6,\y+3.5);
    \foreach \x/\y in {-8/7,8/7,-8/-1,8/-1} \node[font=\figfont\scriptsize] at (\x,\y) {核心};
    \draw[fill=white, draw=rulegray] (-14,-12) rectangle (14,-6);
    \node[font=\figfont\scriptsize] at (0,-9) {大容量缓存};
    \node[anchor=north] at (0,-16) {CPU};
    \node[note, anchor=north] at (0,-21) {少量复杂核心};
  \end{scope}
  % GPU：大量简单核心
  \begin{scope}[shift={(0,0)}]
    \draw[figline, fill=white] (-17,-14) rectangle (17,14);
    \foreach \i in {0,...,7} \foreach \j in {0,...,5}
      \draw[fill=figfill, line width=0.3pt] (-14.8+\i*3.8,-11+\j*3.8) rectangle ++(3,3);
    \node[anchor=north] at (0,-16) {GPU};
    \node[note, anchor=north] at (0,-21) {大量简单并行单元};
  \end{scope}
  % NPU：矩阵运算阵列
  \begin{scope}[shift={(48,0)}]
    \draw[figline, fill=white] (-17,-14) rectangle (17,14);
    \foreach \i in {0,...,5} \foreach \j in {0,...,5}
      \draw[fill=accentpale, draw=accent, line width=0.3pt] (-12+\i*4.2,-12+\j*4.2) rectangle ++(3.4,3.4);
    \foreach \j in {0,...,5} \draw[->, accent, line width=0.4pt] (-16,-10.3+\j*4.2) -- (-12.4,-10.3+\j*4.2);
    \node[anchor=north] at (0,-16) {NPU / ASIC};
    \node[note, anchor=north] at (0,-21) {专用矩阵运算阵列};
  \end{scope}
  % 光谱箭头
  \draw[<->, line width=0.7pt] (-65,-30) -- (65,-30);
  \node[note, anchor=north west] at (-65,-31) {通用性、灵活性更强};
  \node[note, anchor=north east] at (65,-31) {专用计算效率更高};
\end{tikzpicture}
\caption{三类计算芯片的结构取舍（示意）。从左到右，芯片的通用性逐渐降低，矩阵运算的吞吐与能效逐渐提高}
\label{fig:chips}
\end{figure}
}
\Needspace{70mm}\setcounter{figure}{2}\noindent\textbf{图 4-3\quad 版本 3：字体 + 黑灰 + 直角}
{\fontbook\monobook\begin{figure}[htbp]
\centering
\begin{tikzpicture}[fig, x=1mm, y=1mm]
  % CPU：少量复杂核心 + 大缓存
  \begin{scope}[shift={(-48,0)}]
    \draw[figline, fill=white] (-17,-14) rectangle (17,14);
    \foreach \x/\y in {-8/7,8/7,-8/-1,8/-1}
      \draw[fill=figfill] (\x-6,\y-3.5) rectangle (\x+6,\y+3.5);
    \foreach \x/\y in {-8/7,8/7,-8/-1,8/-1} \node[font=\figfont\scriptsize] at (\x,\y) {核心};
    \draw[fill=white, draw=rulegray] (-14,-12) rectangle (14,-6);
    \node[font=\figfont\scriptsize] at (0,-9) {大容量缓存};
    \node[anchor=north] at (0,-16) {CPU};
    \node[note, anchor=north] at (0,-21) {少量复杂核心};
  \end{scope}
  % GPU：大量简单核心
  \begin{scope}[shift={(0,0)}]
    \draw[figline, fill=white] (-17,-14) rectangle (17,14);
    \foreach \i in {0,...,7} \foreach \j in {0,...,5}
      \draw[fill=figfill, line width=0.3pt] (-14.8+\i*3.8,-11+\j*3.8) rectangle ++(3,3);
    \node[anchor=north] at (0,-16) {GPU};
    \node[note, anchor=north] at (0,-21) {大量简单并行单元};
  \end{scope}
  % NPU：矩阵运算阵列
  \begin{scope}[shift={(48,0)}]
    \draw[figline, fill=white] (-17,-14) rectangle (17,14);
    \foreach \i in {0,...,5} \foreach \j in {0,...,5}
      \draw[fill=accentpale, draw=accent, line width=0.3pt] (-12+\i*4.2,-12+\j*4.2) rectangle ++(3.4,3.4);
    \foreach \j in {0,...,5} \draw[->, accent, line width=0.4pt] (-16,-10.3+\j*4.2) -- (-12.4,-10.3+\j*4.2);
    \node[anchor=north] at (0,-16) {NPU / ASIC};
    \node[note, anchor=north] at (0,-21) {专用矩阵运算阵列};
  \end{scope}
  % 光谱箭头
  \draw[<->, line width=0.7pt] (-65,-30) -- (65,-30);
  \node[note, anchor=north west] at (-65,-31) {通用性、灵活性更强};
  \node[note, anchor=north east] at (65,-31) {专用计算效率更高};
\end{tikzpicture}
\caption{三类计算芯片的结构取舍（示意）。从左到右，芯片的通用性逐渐降低，矩阵运算的吞吐与能效逐渐提高}
\label{figsq:chips}
\end{figure}
}

\clearpage
\Needspace{70mm}\setcounter{figure}{6}\noindent\textbf{图 4-7\quad 版本 1：现行版}
\begin{figure}[htbp]
\centering
\begin{tikzpicture}[fig, x=1mm, y=1mm]
  % 三个面板共用同一网格：设备框高 36 mm（y = -18 ~ 18），层的行高 5.5 mm，数据框在 y = -12
  \tikzset{lay/.style={draw=figline, fill=figfill, minimum width=14mm, minimum height=4.2mm, inner sep=0pt, font=\figfont\scriptsize},
           half/.style={draw=accent, fill=accentpale, minimum width=8mm, minimum height=4.2mm, inner sep=0pt, font=\figfont\scriptsize},
           dat/.style={draw=figline, fill=white, rounded corners=1pt, minimum height=4.2mm, inner sep=0pt, font=\figfont\scriptsize},
           dev/.style={draw=rulegray, rounded corners=1.5pt, minimum width=18mm, minimum height=36mm},
           gpu/.style={font=\figfont\scriptsize, anchor=south}}
  \def\row#1{(13.5-#1*5.5)}   % 第 #1 层的 y 坐标：8, 2.5, -3, -8.5
  % ---------- (a) 数据并行 ----------
  \begin{scope}[shift={(-48,0)}]
    \foreach \g/\x in {1/-12, 2/12} {
      \node[dev] (a\g) at (\x,0) {};
      \node[gpu] at (a\g.north) {GPU\,\g};
      \foreach \l in {1,...,4} \node[lay] at (\x,{13.5-\l*5.5}) {第\,\l\,层};
      \node[dat, minimum width=14mm] at (\x,-14) {数据 \g};
    }
    \draw[<->, accent, line width=0.9pt] (-3,0) -- (3,0);
    \node[note, font=\figfont\scriptsize, text=accent] at (0,-23) {梯度 All-Reduce};
    \node[font=\figfont\footnotesize, anchor=north] at (0,-27) {(a) 数据并行：复制模型};
  \end{scope}
  % ---------- (b) 张量并行 ----------
  \begin{scope}[shift={(0,0)}]
    \node[dev, minimum width=30mm] (tp) at (0,0) {};
    \node[gpu] at (-7.5,18) {GPU\,1};
    \node[gpu] at ( 7.5,18) {GPU\,2};
    \draw[densely dashed, rulegray] (0,16) -- (0,-10.5);
    \foreach \l in {1,...,4} {
      \node[half] at (-7.5,{13.5-\l*5.5}) {\l\,a};
      \node[half] at ( 7.5,{13.5-\l*5.5}) {\l\,b};
      \draw[<->, warn, line width=0.6pt] (-3.1,{13.5-\l*5.5}) -- (3.1,{13.5-\l*5.5});
    }
    \node[dat, minimum width=22mm] at (0,-14) {同一批数据};
    \node[note, font=\figfont\scriptsize, text=warn] at (0,-23) {每层都要交换结果};
    \node[font=\figfont\footnotesize, anchor=north] at (0,-27) {(b) 张量并行：切分层内矩阵};
  \end{scope}
  % ---------- (c) 流水线并行 ----------
  \begin{scope}[shift={(48,0)}]
    \foreach \g/\x in {1/-12, 2/12} {
      \node[dev] (p\g) at (\x,0) {};
      \node[gpu] at (p\g.north) {GPU\,\g};
    }
    \node[lay] at (-12,8)   {第\,1\,层};
    \node[lay] at (-12,2.5) {第\,2\,层};
    \node[lay] at ( 12,2.5) {第\,3\,层};
    \node[lay] at ( 12,-3)  {第\,4\,层};
    \node[dat, minimum width=14mm] at (-12,-14) {数据};
    \draw[->, accent, line width=0.9pt] (-5,2.5) -- (5,2.5);
    \node[note, font=\figfont\scriptsize, text=accent] at (0,-23) {阶段间传递激活};
    \node[font=\figfont\footnotesize, anchor=north] at (0,-27) {(c) 流水线并行：按层切分};
  \end{scope}
\end{tikzpicture}
\caption{三种基本并行方式的切分对象（以 2 张 GPU、4 层模型为例）。红色箭头表示频繁通信，应放在高带宽链路上}
\label{fig:parallel}
\end{figure}

\Needspace{70mm}\setcounter{figure}{6}\noindent\textbf{图 4-7\quad 版本 2：只换字体}
{\fontbook\begin{figure}[htbp]
\centering
\begin{tikzpicture}[fig, x=1mm, y=1mm]
  % 三个面板共用同一网格：设备框高 36 mm（y = -18 ~ 18），层的行高 5.5 mm，数据框在 y = -12
  \tikzset{lay/.style={draw=figline, fill=figfill, minimum width=14mm, minimum height=4.2mm, inner sep=0pt, font=\figfont\scriptsize},
           half/.style={draw=accent, fill=accentpale, minimum width=8mm, minimum height=4.2mm, inner sep=0pt, font=\figfont\scriptsize},
           dat/.style={draw=figline, fill=white, rounded corners=1pt, minimum height=4.2mm, inner sep=0pt, font=\figfont\scriptsize},
           dev/.style={draw=rulegray, rounded corners=1.5pt, minimum width=18mm, minimum height=36mm},
           gpu/.style={font=\figfont\scriptsize, anchor=south}}
  \def\row#1{(13.5-#1*5.5)}   % 第 #1 层的 y 坐标：8, 2.5, -3, -8.5
  % ---------- (a) 数据并行 ----------
  \begin{scope}[shift={(-48,0)}]
    \foreach \g/\x in {1/-12, 2/12} {
      \node[dev] (a\g) at (\x,0) {};
      \node[gpu] at (a\g.north) {GPU\,\g};
      \foreach \l in {1,...,4} \node[lay] at (\x,{13.5-\l*5.5}) {第\,\l\,层};
      \node[dat, minimum width=14mm] at (\x,-14) {数据 \g};
    }
    \draw[<->, accent, line width=0.9pt] (-3,0) -- (3,0);
    \node[note, font=\figfont\scriptsize, text=accent] at (0,-23) {梯度 All-Reduce};
    \node[font=\figfont\footnotesize, anchor=north] at (0,-27) {(a) 数据并行：复制模型};
  \end{scope}
  % ---------- (b) 张量并行 ----------
  \begin{scope}[shift={(0,0)}]
    \node[dev, minimum width=30mm] (tp) at (0,0) {};
    \node[gpu] at (-7.5,18) {GPU\,1};
    \node[gpu] at ( 7.5,18) {GPU\,2};
    \draw[densely dashed, rulegray] (0,16) -- (0,-10.5);
    \foreach \l in {1,...,4} {
      \node[half] at (-7.5,{13.5-\l*5.5}) {\l\,a};
      \node[half] at ( 7.5,{13.5-\l*5.5}) {\l\,b};
      \draw[<->, warn, line width=0.6pt] (-3.1,{13.5-\l*5.5}) -- (3.1,{13.5-\l*5.5});
    }
    \node[dat, minimum width=22mm] at (0,-14) {同一批数据};
    \node[note, font=\figfont\scriptsize, text=warn] at (0,-23) {每层都要交换结果};
    \node[font=\figfont\footnotesize, anchor=north] at (0,-27) {(b) 张量并行：切分层内矩阵};
  \end{scope}
  % ---------- (c) 流水线并行 ----------
  \begin{scope}[shift={(48,0)}]
    \foreach \g/\x in {1/-12, 2/12} {
      \node[dev] (p\g) at (\x,0) {};
      \node[gpu] at (p\g.north) {GPU\,\g};
    }
    \node[lay] at (-12,8)   {第\,1\,层};
    \node[lay] at (-12,2.5) {第\,2\,层};
    \node[lay] at ( 12,2.5) {第\,3\,层};
    \node[lay] at ( 12,-3)  {第\,4\,层};
    \node[dat, minimum width=14mm] at (-12,-14) {数据};
    \draw[->, accent, line width=0.9pt] (-5,2.5) -- (5,2.5);
    \node[note, font=\figfont\scriptsize, text=accent] at (0,-23) {阶段间传递激活};
    \node[font=\figfont\footnotesize, anchor=north] at (0,-27) {(c) 流水线并行：按层切分};
  \end{scope}
\end{tikzpicture}
\caption{三种基本并行方式的切分对象（以 2 张 GPU、4 层模型为例）。红色箭头表示频繁通信，应放在高带宽链路上}
\label{fig:parallel}
\end{figure}
}
\Needspace{70mm}\setcounter{figure}{6}\noindent\textbf{图 4-7\quad 版本 3：字体 + 黑灰 + 直角}
{\fontbook\monobook\begin{figure}[htbp]
\centering
\begin{tikzpicture}[fig, x=1mm, y=1mm]
  % 三个面板共用同一网格：设备框高 36 mm（y = -18 ~ 18），层的行高 5.5 mm，数据框在 y = -12
  \tikzset{lay/.style={draw=figline, fill=figfill, minimum width=14mm, minimum height=4.2mm, inner sep=0pt, font=\figfont\scriptsize},
           half/.style={draw=accent, fill=accentpale, minimum width=8mm, minimum height=4.2mm, inner sep=0pt, font=\figfont\scriptsize},
           dat/.style={draw=figline, fill=white, minimum height=4.2mm, inner sep=0pt, font=\figfont\scriptsize},
           dev/.style={draw=rulegray, minimum width=18mm, minimum height=36mm},
           gpu/.style={font=\figfont\scriptsize, anchor=south}}
  \def\row#1{(13.5-#1*5.5)}   % 第 #1 层的 y 坐标：8, 2.5, -3, -8.5
  % ---------- (a) 数据并行 ----------
  \begin{scope}[shift={(-48,0)}]
    \foreach \g/\x in {1/-12, 2/12} {
      \node[dev] (a\g) at (\x,0) {};
      \node[gpu] at (a\g.north) {GPU\,\g};
      \foreach \l in {1,...,4} \node[lay] at (\x,{13.5-\l*5.5}) {第\,\l\,层};
      \node[dat, minimum width=14mm] at (\x,-14) {数据 \g};
    }
    \draw[<->, accent, line width=0.9pt] (-3,0) -- (3,0);
    \node[note, font=\figfont\scriptsize, text=accent] at (0,-23) {梯度 All-Reduce};
    \node[font=\figfont\footnotesize, anchor=north] at (0,-27) {(a) 数据并行：复制模型};
  \end{scope}
  % ---------- (b) 张量并行 ----------
  \begin{scope}[shift={(0,0)}]
    \node[dev, minimum width=30mm] (tp) at (0,0) {};
    \node[gpu] at (-7.5,18) {GPU\,1};
    \node[gpu] at ( 7.5,18) {GPU\,2};
    \draw[densely dashed, rulegray] (0,16) -- (0,-10.5);
    \foreach \l in {1,...,4} {
      \node[half] at (-7.5,{13.5-\l*5.5}) {\l\,a};
      \node[half] at ( 7.5,{13.5-\l*5.5}) {\l\,b};
      \draw[<->, warn, line width=0.6pt] (-3.1,{13.5-\l*5.5}) -- (3.1,{13.5-\l*5.5});
    }
    \node[dat, minimum width=22mm] at (0,-14) {同一批数据};
    \node[note, font=\figfont\scriptsize, text=warn] at (0,-23) {每层都要交换结果};
    \node[font=\figfont\footnotesize, anchor=north] at (0,-27) {(b) 张量并行：切分层内矩阵};
  \end{scope}
  % ---------- (c) 流水线并行 ----------
  \begin{scope}[shift={(48,0)}]
    \foreach \g/\x in {1/-12, 2/12} {
      \node[dev] (p\g) at (\x,0) {};
      \node[gpu] at (p\g.north) {GPU\,\g};
    }
    \node[lay] at (-12,8)   {第\,1\,层};
    \node[lay] at (-12,2.5) {第\,2\,层};
    \node[lay] at ( 12,2.5) {第\,3\,层};
    \node[lay] at ( 12,-3)  {第\,4\,层};
    \node[dat, minimum width=14mm] at (-12,-14) {数据};
    \draw[->, accent, line width=0.9pt] (-5,2.5) -- (5,2.5);
    \node[note, font=\figfont\scriptsize, text=accent] at (0,-23) {阶段间传递激活};
    \node[font=\figfont\footnotesize, anchor=north] at (0,-27) {(c) 流水线并行：按层切分};
  \end{scope}
\end{tikzpicture}
\caption{三种基本并行方式的切分对象（以 2 张 GPU、4 层模型为例）。粗双向箭头表示频繁通信，应放在高带宽链路上}
\label{figsq:parallel}
\end{figure}
}
\end{document}
